% sec9_second.tex: the 1/n term for pure powers with p>1 (Theorem thm:onen).

\section{\texorpdfstring{The $1/n$ term for $p>1$}{The 1/n term for p>1}}\label{sec:second}

This section proves Theorem~\ref{thm:onen} by expanding the error matrix $\mathsf R$ of \S\ref{sec:classical-small}
to second order in $1/n$. The standing assumptions and the notation of Section~\ref{sec:classical} are in force: $p>1$,
$s>1$, $\ell\equiv1$, and $\alpha>-1$ is arbitrary.

\emph{Outline.} The global parametrix $P^\infty_0$ does not depend on $n$ or on the weight, and it contributes exactly
$c^2$ and $2c$ to the two extraction functionals. So the $1/n$ terms of $b_n$ and $a_n$ come only from the first
corrections of the two local models, the Bessel model at $0$ and the Airy model at $4c$, on the two circles
(Lemmas~\ref{lem:model-corrections} and~\ref{lem:circle-jumps}). These corrections depend on $n$ only through the band
factor $H_n$, and $H_n$ is constant up to $O(n^{-p})$ because the two $1/(12n)$ Binet terms cancel
(Lemma~\ref{lem:bandfactor}); so at order $1/n$ the corrections are those of the fixed model weight $x^\alpha(4c-x)$,
whose local data are explicit (Lemma~\ref{lem:local-data}). The corrections are meromorphic in the two discs, and the
residues enter the coefficients at infinity (Proposition~\ref{prop:second}). The residue at the hard edge is a multiple of
$\bigl(\begin{smallmatrix}1&i\\i&-1\end{smallmatrix}\bigr)$, which both extraction functionals annihilate; so the $1/n$ terms
come entirely from the Airy correction, through two combinations of its Laurent coefficients that we compute by hand
(Lemma~\ref{lem:laurent}).

Throughout, $\mathsf T:=\frac1{\sqrt2}\bigl(\begin{smallmatrix}1&1\\i&-i\end{smallmatrix}\bigr)$, so that
$P^\infty_0=\mathsf Ta^{\sigma_3}\mathsf T^{-1}$ with the function $a(z)=((z-4c)/z)^{1/4}$ of~\eqref{eq:Pinf0} (it is never
indexed, and is not to be confused with the coefficients $a_n$). For a matrix $Z$ with $Z_{22}=-Z_{11}$, a direct
computation gives
\begin{equation}\label{eq:conj-T}
  P^\infty_0Z(P^\infty_0)^{-1}=\mathsf T\begin{pmatrix}Y_{11}&a^2Y_{12}\\a^{-2}Y_{21}&-Y_{11}\end{pmatrix}\mathsf T^{-1},\qquad
  \begin{aligned}Y_{11}&=\tfrac i2(Z_{12}-Z_{21}),\\ Y_{12}&=Z_{11}-\tfrac i2(Z_{12}+Z_{21}),\\ Y_{21}&=Z_{11}+\tfrac i2(Z_{12}+Z_{21}),\end{aligned}
\end{equation}
and, for any numbers $y_{11},y_{12},y_{21}$,
\begin{equation}\label{eq:T-back}
  \mathsf T\begin{pmatrix}y_{11}&y_{12}\\y_{21}&-y_{11}\end{pmatrix}\mathsf T^{-1}=
  \begin{pmatrix}\frac12(y_{12}+y_{21})&\frac i2(y_{12}-y_{21}-2y_{11})\\[1mm]\frac i2(y_{12}-y_{21}+2y_{11})&-\frac12(y_{12}+y_{21})\end{pmatrix}.
\end{equation}

\subsection{The first corrections}\label{sec:second-corrections}

\begin{lemma}[First corrections of the two models]\label{lem:model-corrections}
\begin{enumerate}
\item[(a)] Let $\alpha>-1$ and
  \[
    \mathsf Z_{\mathrm B}:=\begin{pmatrix}-(1+4\alpha^2)&2i\\2i&1+4\alpha^2\end{pmatrix}.
  \]
  Uniformly as $\zeta\to\infty$ off the rays of $\Psi_\alpha$,
  \[
    \Psi_\alpha(\zeta)=(2\pi\zeta^{1/2})^{-\sigma_3/2}M\Bigl[I+\frac{\sigma_3\mathsf Z_{\mathrm B}\sigma_3}{16\zeta^{1/2}}
    +O(\zeta^{-1})\Bigr]e^{2\zeta^{1/2}\sigma_3},
  \]
  equivalently $\Phi_\alpha(\zeta)=(2\pi\zeta^{1/2})^{-\sigma_3/2}M^{-1}\bigl[I+\mathsf Z_{\mathrm B}/(16\zeta^{1/2})+O(\zeta^{-1})\bigr]$.
\item[(b)] Let $\mathsf Z_{\mathrm{Ai}}:=\bigl(\begin{smallmatrix}-1/72&-i/12\\-i/12&1/72\end{smallmatrix}\bigr)$ and
  $\xi:=\frac23\zeta^{3/2}$. Uniformly for $\zeta\in\mathcal S_j$, $|\zeta|\ge1$, the Airy model~\eqref{eq:airy-model} satisfies
  \[
    \Psi^j(\zeta)=\zeta^{-\sigma_3/4}M\bigl(I+\mathsf Z_{\mathrm{Ai}}/\xi+O(\xi^{-2})\bigr)e^{\xi\sigma_3}\qquad(j=\mathrm I,\dots,\mathrm{IV}).
  \]
\end{enumerate}
\end{lemma}

Part~(a) uses \cite[(6.23)--(6.25)]{KMVV} and the large-argument expansions (10.40.1)--(10.40.4) with the coefficients
(10.17.1), (10.17.8), and (10.17.5), (10.17.6), (10.17.11), (10.17.12) for the Hankel functions
(\S\ref{sec:tool-dlmf}(d)); these hold for every fixed order, so no restriction on $\alpha$ arises. Part~(b) uses
(9.7.2), (9.7.5) and (9.7.6) (\S\ref{sec:tool-dlmf}(b)). Both are the computations behind \cite[(6.28)]{KMVV}
and~\eqref{eq:airy-infty}, carried one term further.

\begin{proof}
(a) In $|\arg\zeta|<2\pi/3$, \cite[(6.23)]{KMVV} expresses $\Psi_\alpha$ through $I_\alpha$, $K_\alpha$ and their
derivatives at $w=2\zeta^{1/2}$, and $|\mathrm{ph}\,w|<\pi/3$, inside the sectors of (10.40.1)--(10.40.4). With the first
coefficients $a_1(\alpha)=(4\alpha^2-1)/8$ and $b_1(\alpha)=(4\alpha^2+3)/8$ these give, entrywise with relative errors
$O(w^{-2})$, $\Psi_\alpha e^{-w\sigma_3}=(2\pi\zeta^{1/2})^{-\sigma_3/2}\bigl[M+w^{-1}\frac1{\sqrt2}
\bigl(\begin{smallmatrix}-a_1&ia_1\\-ib_1&b_1\end{smallmatrix}\bigr)\bigr]$. Since
\[
  M^{-1}\frac1{\sqrt2}\begin{pmatrix}-a_1&ia_1\\-ib_1&b_1\end{pmatrix}=\frac12\begin{pmatrix}-(a_1+b_1)&i(a_1-b_1)\\i(a_1-b_1)&a_1+b_1\end{pmatrix},
  \qquad a_1+b_1=\frac{4\alpha^2+1}4,\quad a_1-b_1=-\frac12,
\]
the correction is $\sigma_3\mathsf Z_{\mathrm B}\sigma_3/(16\zeta^{1/2})$. In the sectors $2\pi/3<\pm\arg\zeta<\pi$,
\cite[(6.24)--(6.25)]{KMVV} use $H^{(1,2)}_\alpha$ and their derivatives at $2(-\zeta)^{1/2}$, whose phase lies in
$(-\pi/6,\pi/6)$, inside the sectors of (10.17.5)--(10.17.6) and (10.17.11)--(10.17.12); the same computation with the first two
terms gives the same formula. The form for $\Phi_\alpha$ follows from $\sigma_3M\sigma_3=M^{-1}$.

(b) In $\mathcal S_{\mathrm I}$ the columns of $\Psi^{\mathrm I}$ are $k_{\mathrm{Ai}}e^{-i\pi\sigma_3/4}y_+$ and
$k_{\mathrm{Ai}}e^{-i\pi\sigma_3/4}\mathsf qy_0$. By (9.7.5)--(9.7.6) with $u_1=\frac5{72}$, $v_1=-\frac7{72}$ from (9.7.2), applied
at $\zeta$ and at $\mathsf q^{-1}\zeta$ (where $(\mathsf q^{-1}\zeta)^{1/4}=e^{-i\pi/6}\zeta^{1/4}$ and
$\frac23(\mathsf q^{-1}\zeta)^{3/2}=-\xi$), with relative errors $O(\xi^{-2})$,
\[
  y_+=\frac{e^\xi e^{i\pi/6}}{2\sqrt\pi}\zeta^{-\sigma_3/4}\begin{pmatrix}1+u_1/\xi\\1+v_1/\xi\end{pmatrix},\qquad
  y_0=\frac{e^{-\xi}}{2\sqrt\pi}\zeta^{-\sigma_3/4}\begin{pmatrix}1-u_1/\xi\\-(1-v_1/\xi)\end{pmatrix}.
\]
So $\Psi^{\mathrm I}=\zeta^{-\sigma_3/4}(N_0+N_1/\xi+O(\xi^{-2}))e^{\xi\sigma_3}$ with $N_0=M$ (this is~\eqref{eq:airy-infty}) and
$M^{-1}N_1=\bigl(\begin{smallmatrix}(u_1+v_1)/2&-i(u_1-v_1)/2\\-i(u_1-v_1)/2&-(u_1+v_1)/2\end{smallmatrix}\bigr)=\mathsf Z_{\mathrm{Ai}}$,
using $\mathsf qe^{-i\pi/6}=i$. In each of the other sectors the two columns of $\Psi^j$ are evaluated at points with the
same leading asymptotics as in $\mathcal S_{\mathrm I}$ (this is the computation behind~\eqref{eq:airy-infty}; for
instance, in $\mathcal S_{\mathrm{II}}$, $-\mathsf q^{-1}y_-$ has the same two-term expansion as $\mathsf qy_0$, since
$(\mathsf q\zeta)^{1/4}=e^{-i\pi/3}\zeta^{1/4}$ and $(\mathsf q\zeta)^{3/2}=\zeta^{3/2}$ there), so the correction is the same.
\end{proof}

Let $f^{1/2}$ be the principal root; it is analytic and non-zero on $\partial U_0$. Put, on $\partial U_0$ and $\partial U_r$
respectively,
\begin{equation}\label{eq:Delta}
  \Delta_0:=\frac1{16f^{1/2}}\,P^\infty_0\rho_n^{\sigma_3}\mathsf Z_{\mathrm B}\rho_n^{-\sigma_3}(P^\infty_0)^{-1},\qquad
  \Delta_r:=-\frac1{\tvphi}\,P^\infty_0(-\omega)^{-\sigma_3/2}\mathsf Z_{\mathrm{Ai}}(-\omega)^{\sigma_3/2}(P^\infty_0)^{-1},
\end{equation}
with $\rho_n$ of~\eqref{eq:Wtilde} and the root~\eqref{eq:omega-root}. Here $\Delta_0$ depends on $\rho_n$ only
through $\rho_n^{\pm2}=(e^{\mp i\pi\alpha}/\omega)^{\pm1}$, and $\Delta_r$ only through $(-\omega)^{\pm1}$. Let
$\Delta_0^{\mathrm{mod}}$ and $\Delta_r^{\mathrm{mod}}$ be the same expressions with $\omega$ replaced by $\omega^{\mathrm{mod}}$ of
\eqref{eq:omega-mod}, on the whole punctured discs $\mathbb D(0,2\delta)\setminus\{0\}$ and $\mathbb D(4c,2\eps)\setminus\{4c\}$ (off
$[0,4c]$ to begin with).

\begin{lemma}[Jumps on the two circles]\label{lem:circle-jumps}
Uniformly on the circles, with boundary values where they cross the contour,
\[
  \check P_0(P^\infty_0)^{-1}=I+\frac{\Delta_0}n+O(n^{-2})\ \text{on }\partial U_0,\qquad
  \hat P_r(P^\infty_0)^{-1}=I+\frac{\Delta_r}n+O(n^{-2})\ \text{on }\partial U_r,
\]
and $\Delta_0-\Delta_0^{\mathrm{mod}}=O(n^{-p})$, $\Delta_r-\Delta_r^{\mathrm{mod}}=O(n^{-p})$ there.
\end{lemma}

\begin{proof}
\emph{Bessel.} On $\partial U_0$, $\zeta=n^2f$ and $\zeta^{1/2}=nf^{1/2}$, with $|\zeta|\asymp n^2$. By~\eqref{eq:Pcheck} and
Lemma~\ref{lem:model-corrections}(a), since $f^{\sigma_3/4}(2\pi n)^{\sigma_3/2}(2\pi nf^{1/2})^{-\sigma_3/2}=I$,
\[
  \check P_0=P^\infty_0\rho_n^{\sigma_3}MM^{-1}\Bigl[I+\frac{\mathsf Z_{\mathrm B}}{16nf^{1/2}}+O(n^{-2})\Bigr]\rho_n^{-\sigma_3},
\]
which is the first claim, because $P^\infty_0$, $(P^\infty_0)^{-1}$ and $\rho_n^{\pm1}$ are bounded on $\partial U_0$
uniformly in $n$ (Lemma~\ref{lem:classical-param}(i)).

\emph{Airy.} On $\partial U_r$, by~\eqref{eq:airy}, $\hat P_r=P^\infty_0(-\omega)^{-\sigma_3/2}M^{-1}f_r^{\sigma_3/4}n^{\sigma_3/6}
\Psi_r(\zeta)((-\omega)^{1/2}e^{n\tvphi})^{\sigma_3}$ with $\zeta=n^{2/3}f_r$, so $\zeta^{-\sigma_3/4}=n^{-\sigma_3/6}f_r^{-\sigma_3/4}$
and $\xi=\frac23nf_r^{3/2}=-n\tvphi$ (Lemma~\ref{lem:signs}(iv)). Lemma~\ref{lem:model-corrections}(b) and
$e^{\xi\sigma_3}e^{n\tvphi\sigma_3}=I$ give
\[
  \hat P_r=P^\infty_0(-\omega)^{-\sigma_3/2}\Bigl(I-\frac{\mathsf Z_{\mathrm{Ai}}}{n\tvphi}+O(n^{-2})\Bigr)(-\omega)^{\sigma_3/2},
\]
uniformly on $\partial U_r$ because $|\xi|\asymp n$ there and $P^\infty_0$, $(P^\infty_0)^{-1}$, $(-\omega)^{\pm1/2}$ are bounded
(Lemmas~\ref{lem:szego}(vi)(b) and~\ref{lem:airy}).

\emph{Dependence on $n$.} The functions $f$, $\tvphi$ and $P^\infty_0$ do not depend on $n$. By
Lemma~\ref{lem:bandfactor}(iv), $\omega/\omega^{\mathrm{mod}}=1+O(n^{-p})$ on both circles, and $\omega$, $\omega^{\mathrm{mod}}$
are bounded above and below there; since $\Delta_0$, $\Delta_r$ are smooth functions of $\omega^{\pm1}$ with these bounds,
the last claim follows.
\end{proof}

\subsection{The local data and the Laurent coefficients}\label{sec:second-laurent}

\begin{lemma}[Local data of the model]\label{lem:local-data}
\begin{enumerate}
\item[(i)] $f(z)=f'(0)z\bigl(1+f_2z+O(z^2)\bigr)$ near $0$, with
  \[
    f'(0)=-\frac1{c(p+2)^2},\qquad f_2=\frac{p+2}{12c(3p+2)}.
  \]
\item[(ii)] With $\mathsf u:=(z-4c)^{1/2}$ (principal root), $\tvphi=-\mathsf g_3\mathsf u^3-\mathsf g_5\mathsf u^5+O(\mathsf u^7)$ near
  $4c$, where $\tvphi/\mathsf u^3$ is analytic at $4c$ and
  \[
    \mathsf g_3=\frac1{6pc^{3/2}},\qquad \mathsf g_5=-\frac{7p+4}{240p^2c^{5/2}}.
  \]
\item[(iii)] Near $4c$ and near $0$ (in $\C_\pm$), respectively,
  \[
    (-\omega^{\mathrm{mod}})^{1/2}=\Bigl(\frac{2\sqrt c}{z^{1/2}+\mathsf u}\Bigr)^{1+\alpha},\qquad
    \frac{e^{\mp i\pi\alpha}}{\omega^{\mathrm{mod}}}=\Bigl(\frac{((-z)^{1/2}+(4c-z)^{1/2})^2}{4c}\Bigr)^{1+\alpha};
  \]
  equivalently, with $\mathsf y:=(-z)^{1/2}$,
  \[
    L_r:=(1+\alpha)\operatorname{arsinh}\frac{\mathsf u}{2\sqrt c},\quad L_0:=(1+\alpha)\operatorname{arsinh}\frac{\mathsf y}{2\sqrt c},
    \quad(-\omega^{\mathrm{mod}})^{\mp1}=e^{\pm2L_r},\quad\frac{e^{\mp i\pi\alpha}}{\omega^{\mathrm{mod}}}=e^{2L_0}.
  \]
\end{enumerate}
\end{lemma}

The lemma uses only the equilibrium measure (Proposition~\ref{prop:equilibrium}, from \cite[Theorem~1.3]{Paper1}),
through the formulas for $f$ and $\tvphi$ in the proof of Lemma~\ref{lem:signs}, and Lemma~\ref{lem:szego}(i).

\begin{proof}
Put $a':=\frac1p+\frac12$; recall $(4c)^{-1/p}=2I(p)/\pi$ and $B(a',\frac12)=2I(p)$ by~\eqref{eq:cp} and (5.12.1).

(i) By Proposition~\ref{prop:equilibrium}, on $(0,4c)$, $\psi^*(x)=\lambda_p(x)\mathcal B_{x/4c}(a',\frac12)$, and
$\mathcal B_y(a',\frac12)=\frac{y^{a'}}{a'B(a',\frac12)}\bigl(1+\frac{a'y}{2(a'+1)}+O(y^2)\bigr)$. Hence
$\psi^*(x)=C_*x^{-1/2}(1+\tau_*x+O(x^2))$ with $C_*=(4c)^{-1/2}/(\pi(1+p/2))$ (Proposition~\ref{prop:equilibrium}(iii)) and
$\tau_*=a'/(8c(a'+1))=(p+2)/(8c(3p+2))$. So $\mu^*((0,x])=x^{1/2}\mathfrak a(x)$ with $\mathfrak a(z)=2C_*(1+\tau_*z/3+O(z^2))$,
and $f=-\frac{\pi^2}4z\mathfrak a^2$ (proof of Lemma~\ref{lem:signs}(iii)) gives $f'(0)=-\pi^2C_*^2=-1/(c(p+2)^2)$ and
$f_2=2\tau_*/3$.

(ii) By the proof of Lemma~\ref{lem:signs}(iv), $\tvphi_+(x)=i\pi\int_x^{4c}(\lambda_p-\psi^*)$ on the band. With
$y=4c(1-e)$, $1-\mathcal B_{1-e}(a',\frac12)=\frac1{B(a',1/2)}\bigl(2e^{1/2}-\frac23(a'-1)e^{3/2}+O(e^{5/2})\bigr)$ and
$\lambda_p(y)=\frac1p(4c)^{-1-1/p}(1+(1+\frac1p)e+O(e^2))$; integrating, with $v:=4c-x$,
\[
  \int_x^{4c}(\lambda_p-\psi^*)=\frac1{\pi p}\Bigl[\frac43\Bigl(\frac v{4c}\Bigr)^{3/2}+\frac25\Bigl(\frac73+\frac4{3p}\Bigr)\Bigl(\frac v{4c}\Bigr)^{5/2}
  +O(v^{7/2})\Bigr].
\]
On the band, approached from $\C_+$, $\mathsf u=iv^{1/2}$, so $-\mathsf g_3\mathsf u^3-\mathsf g_5\mathsf u^5=i\mathsf g_3v^{3/2}-i\mathsf g_5v^{5/2}$;
comparing coefficients gives the values. The quotient $\tvphi/\mathsf u^3$ is analytic at $4c$ because
$\tvphi=-\frac23f_r^{3/2}$ with $f_r$ conformal and $f_r(4c)=0$ (Lemma~\ref{lem:signs}(iv)), so the band values determine
the expansion.

(iii) By Lemma~\ref{lem:szego}(i), $\omega^{\mathrm{mod}}=(4cz/\Upsilon)^\alpha z^{-\alpha}\cdot(-4c/\Upsilon)$. Near $4c$ all
arguments are close to $0$, so $-\omega^{\mathrm{mod}}=(4c/\Upsilon)^{1+\alpha}$ and, as in~\eqref{eq:omega-root}, its root is the
stated one; $z^{1/2}=2\sqrt c(1+\mathsf u^2/(4c))^{1/2}$ gives the form with $\operatorname{arsinh}$. Near $0$, in $\C_\pm$,
$z^{1/2}=\pm i(-z)^{1/2}$ and $(z-4c)^{1/2}=\pm i(4c-z)^{1/2}$, so $\Upsilon=-\hat\Upsilon$ with
$\hat\Upsilon:=((-z)^{1/2}+(4c-z)^{1/2})^2$; and $(4cz/\Upsilon)^\alpha=(-z)^\alpha(4c/\hat\Upsilon)^\alpha$, $z^{-\alpha}=(-z)^{-\alpha}e^{\mp
i\pi\alpha}$. Hence $\omega^{\mathrm{mod}}=e^{\mp i\pi\alpha}(4c/\hat\Upsilon)^{1+\alpha}$.
\end{proof}

\begin{lemma}[Laurent coefficients]\label{lem:laurent}
$\Delta_0^{\mathrm{mod}}$ is meromorphic in $\mathbb D(0,2\delta)$ with at most a simple pole at $0$, and $\Delta_r^{\mathrm{mod}}$ is
meromorphic in $\mathbb D(4c,2\eps)$ with at most a double pole at $4c$:
\[
  \Delta_0^{\mathrm{mod}}(z)=\frac{\mathfrak A_0}z+O(1),\qquad
  \Delta_r^{\mathrm{mod}}(z)=\frac{\mathfrak B_r}{(z-4c)^2}+\frac{\mathfrak A_r}{z-4c}+O(1).
\]
Here
\[
  \mathfrak A_0=\frac{c(p+2)(4\alpha^2-1)}{16}\begin{pmatrix}1&i\\i&-1\end{pmatrix},\qquad
  \mathfrak B_r=\frac{5c^2p}{12}\begin{pmatrix}1&-i\\-i&-1\end{pmatrix},
\]
and
\begin{equation}\label{eq:Ar-combinations}
  (\mathfrak A_r)_{21}-(\mathfrak A_r)_{12}=ipc(1+\alpha),\qquad
  (\mathfrak A_r)_{11}-(\mathfrak A_r)_{22}-2i(\mathfrak A_r)_{12}-\frac ic(\mathfrak B_r)_{12}=-pc(\alpha+2).
\end{equation}
The full matrix $\mathfrak A_r$ is given in Lemma~\ref{lem:residues-full}.
\end{lemma}

\begin{proof}
\emph{Meromorphy.} On $(0,2\delta)$, $f^{1/2}_+=-f^{1/2}_-$ (as $f<0$ there and $f$ maps $\C_\pm$ into $\C_\mp$),
$P^\infty_{0,+}=P^\infty_{0,-}J$, and the model value of $\rho^2$ in Lemma~\ref{lem:local-data}(iii) satisfies
$\rho^2_+\rho^2_-=1$; since $J\mathsf Z_{\mathrm B}J^{-1}=-\mathsf Z_{\mathrm B}$ and
$Jd^{-\sigma_3}=d^{\sigma_3}J$ for scalars $d$, $\Delta^{\mathrm{mod}}_{0,+}=\Delta^{\mathrm{mod}}_{0,-}$. On $(-2\delta,0)$ every
factor is analytic. So $\Delta_0^{\mathrm{mod}}$ is analytic in the punctured disc, and $|\Delta_0^{\mathrm{mod}}|\le C|z|^{-1}$
by~\eqref{eq:Pinf0-bound} and $|f|\asymp|z|$; hence $0$ is at most a simple pole. On $(4c-2\eps,4c)$,
$\tvphi_+=-\tvphi_-$ (Lemma~\ref{lem:gfunction}(i)), $(-\omega)^{1/2}_+(-\omega)^{1/2}_-=1$ by~\eqref{eq:omega-root-band} (in
particular for the model), and $J\mathsf Z_{\mathrm{Ai}}J^{-1}=-\mathsf Z_{\mathrm{Ai}}$, so again there is no jump; on
$(4c,4c+2\eps)$ all factors are analytic. Since $|\tvphi|\asymp|z-4c|^{3/2}$ and $|P^\infty_0|\le C|z-4c|^{-1/4}$ near $4c$,
$|\Delta_r^{\mathrm{mod}}|\le C|z-4c|^{-2}$.

\emph{The coefficient $\mathfrak A_0$.} By~\eqref{eq:conj-T} with $Z=\rho^{\sigma_3}\mathsf Z_{\mathrm B}\rho^{-\sigma_3}$ and $\rho^2=e^{2L_0}$,
\[
  Y_{11}=-2\sinh2L_0,\qquad Y_{12}=1-4\alpha^2+(\rho-\rho^{-1})^2,\qquad Y_{21}=-(3+4\alpha^2)-(\rho-\rho^{-1})^2 .
\]
Let $z\to0$ along $(-2\delta,0)$, where $f>0$, $L_0=O(|z|^{1/2})$, $a^2=(4c-z)^{1/2}(-z)^{-1/2}$ and
$f^{1/2}=|f'(0)|^{1/2}(-z)^{1/2}(1+O(z))$. Then $zY_{11}/f^{1/2}\to0$, $za^{-2}Y_{21}/f^{1/2}\to0$, and
$za^2Y_{12}/f^{1/2}\to-2\sqrt c|f'(0)|^{-1/2}(1-4\alpha^2)=-2c(p+2)(1-4\alpha^2)$ by Lemma~\ref{lem:local-data}(i).
Since the pole is at most simple, $\mathfrak A_0=\lim_{z\to0}z\Delta_0^{\mathrm{mod}}(z)$, and~\eqref{eq:T-back} with
$y_{12}=1$, $y_{11}=y_{21}=0$ gives $\mathsf Te_1e_2^{\mathsf T}\mathsf T^{-1}=\frac12\bigl(\begin{smallmatrix}1&i\\i&-1\end{smallmatrix}\bigr)$; so
$\mathfrak A_0=\frac1{16}\cdot(-2c(p+2)(1-4\alpha^2))\cdot\frac12\bigl(\begin{smallmatrix}1&i\\i&-1\end{smallmatrix}\bigr)$.

\emph{The Airy coefficients.} By~\eqref{eq:conj-T} with $Z=(-\omega^{\mathrm{mod}})^{-\sigma_3/2}\mathsf Z_{\mathrm{Ai}}(-\omega^{\mathrm{mod}})^{\sigma_3/2}$,
whose off-diagonal entries are $-\frac i{12}e^{2L_r}$ and $-\frac i{12}e^{-2L_r}$ by Lemma~\ref{lem:local-data}(iii),
\begin{equation}\label{eq:Y-airy}
  Y_{11}=\frac{\sinh2L_r}{12},\qquad Y_{12}=-\frac1{72}-\frac{\cosh2L_r}{12},\qquad Y_{21}=-\frac1{72}+\frac{\cosh2L_r}{12},
\end{equation}
and $\Delta_r^{\mathrm{mod}}=-\tvphi^{-1}\mathsf T\bigl(\begin{smallmatrix}Y_{11}&a^2Y_{12}\\a^{-2}Y_{21}&-Y_{11}\end{smallmatrix}\bigr)\mathsf T^{-1}$.
Take $z\in(4c,4c+2\eps)$, so $\mathsf u>0$, and put $\epsilon:=\mathsf u/(2\sqrt c)$. Then $a^2=\epsilon(1+\epsilon^2)^{-1/2}$ is odd
in $\mathsf u$, $L_r=(1+\alpha)\epsilon+O(\epsilon^3)$, $Y_{11}$ is odd and $Y_{12}$, $Y_{21}$ are even, and $\tvphi^{-1}$ is
$\mathsf u^{-3}$ times an even series. A function meromorphic at $4c$ is an even Laurent series in $\mathsf u$, and its residue
is the coefficient of $\mathsf u^{-2}$. Since the pole is at most double,
$\mathfrak B_r=\lim(z-4c)^2\Delta_r^{\mathrm{mod}}=\mathsf g_3^{-1}\lim\mathsf u\,\mathsf T(\cdots)\mathsf T^{-1}$, and
$\mathsf ua^{-2}Y_{21}\to2\sqrt c\cdot\frac5{72}$ while the other entries times $\mathsf u$ tend to $0$; with
$\mathsf Te_2e_1^{\mathsf T}\mathsf T^{-1}=\frac12\bigl(\begin{smallmatrix}1&-i\\-i&-1\end{smallmatrix}\bigr)$ and $\mathsf g_3^{-1}=6pc^{3/2}$ this
gives $\mathfrak B_r$.

For the first combination in~\eqref{eq:Ar-combinations}, \eqref{eq:T-back} gives
$(\Delta_r^{\mathrm{mod}})_{21}-(\Delta_r^{\mathrm{mod}})_{12}=-2iY_{11}/\tvphi$, and
\[
  \frac{Y_{11}}{\tvphi}=\frac{(1+\alpha)\mathsf u/(12\sqrt c)+O(\mathsf u^3)}{-\mathsf g_3\mathsf u^3(1+O(\mathsf u^2))}
  =-\frac{1+\alpha}{12\sqrt c\,\mathsf g_3}\,\mathsf u^{-2}+O(1);
\]
so the difference has no double pole and its residue is $\frac{i(1+\alpha)}{6\sqrt c\,\mathsf g_3}=ipc(1+\alpha)$. For the second,
$(\mathfrak A_r)_{11}-(\mathfrak A_r)_{22}-2i(\mathfrak A_r)_{12}-\frac ic(\mathfrak B_r)_{12}$ is the residue of
$g:=(\Delta)_{11}-(\Delta)_{22}-2i(\Delta)_{12}-\frac ic(z-4c)(\Delta)_{12}$, $\Delta=\Delta_r^{\mathrm{mod}}$, and~\eqref{eq:T-back} gives
\[
  g=-\frac2{\tvphi}\bigl(a^2Y_{12}-Y_{11}\bigr)-\frac{\mathsf u^2}{2c\,\tvphi}\bigl(a^2Y_{12}-a^{-2}Y_{21}-2Y_{11}\bigr).
\]
In the first term $a^2Y_{12}-Y_{11}=\bigl(-\frac7{144\sqrt c}-\frac{1+\alpha}{12\sqrt c}\bigr)\mathsf u+O(\mathsf u^3)$, an odd series,
which contributes $\frac2{\mathsf g_3}\cdot\bigl(-\frac{12\alpha+19}{144\sqrt c}\bigr)=-\frac{pc(12\alpha+19)}{12}$ to the coefficient of
$\mathsf u^{-2}$. In the second, the bracket is $-\frac{2\sqrt c}{\mathsf u}\cdot\frac5{72}+O(\mathsf u)$, which contributes
$\frac1{2c\,\mathsf g_3}\cdot\bigl(-\frac{10\sqrt c}{72}\bigr)=-\frac{5pc}{12}$. The sum is $-pc(\alpha+2)$.
\end{proof}

The higher Taylor coefficients $f_2$ and $\mathsf g_5$ do not enter~\eqref{eq:Ar-combinations}: only the leading terms of
$a^{\pm2}$, $Y$ and $\tvphi$ were used. They do enter $\mathfrak A_r$ itself and the constant term at $0$
(Lemma~\ref{lem:residues-full}).

\subsection{The second-order small-norm estimate}\label{sec:second-small}

\begin{proposition}[Second-order small norm]\label{prop:second}
Let $p>1$, $s>1$, $\ell\equiv1$. For $n\ge n_0$,
\[
  \mathsf R(z)=I+\frac{\mathsf R^{(1)}(z)}n+O(n^{-2})
\]
uniformly for $\mathrm{dist}(z,\Sigma_{\mathsf R})\ge d>0$, and the coefficients at infinity satisfy $\mathsf R_k=\mathsf R_k^{(1)}/n+O(n^{-2})$,
$k=1,2$, where
\[
  \mathsf R^{(1)}(z)=\frac{\mathfrak A_0}z+\frac{\mathfrak A_r}{z-4c}+\frac{\mathfrak B_r}{(z-4c)^2}\quad(z\notin\overline{U_0}\cup\overline{U_r}),\qquad
  \mathsf R_1^{(1)}=\mathfrak A_0+\mathfrak A_r,\qquad \mathsf R_2^{(1)}=4c\,\mathfrak A_r+\mathfrak B_r .
\]
\end{proposition}

\begin{proof}
By Proposition~\ref{prop:classical}(iii),
\[
  \mathsf R(z)-I=\frac1{2\pi i}\int_{\Sigma_{\mathsf R}}\frac{(J_{\mathsf R}-I)(s)}{s-z}\diff s
  +\frac1{2\pi i}\int_{\Sigma_{\mathsf R}}\frac{(\mathsf R_--I)(J_{\mathsf R}-I)(s)}{s-z}\diff s .
\]
By Cauchy--Schwarz and Proposition~\ref{prop:classical}(ii), (iv) the second integral is at most
$\|\mathsf R_--I\|_{L^2}\|J_{\mathsf R}-I\|_{L^2}/(2\pi d)=O(n^{-2})$. In the first, the arcs inside $U_0$ contribute
$O(n^{-K})$ and the remaining arcs off the circles $O(e^{-n/C})$; by Lemma~\ref{lem:circle-jumps} the circles contribute
\[
  \frac1n\cdot\frac1{2\pi i}\oint_{\mathrm{cw}}\frac{\Delta^{\mathrm{mod}}(s)}{s-z}\diff s+O(n^{-2}+n^{-1-p}),
\]
where $\Delta^{\mathrm{mod}}$ is $\Delta_0^{\mathrm{mod}}$ on $\partial U_0$ and $\Delta_r^{\mathrm{mod}}$ on $\partial U_r$, both circles
oriented clockwise; and $n^{-1-p}\le n^{-2}$ since $p>1$. For $z$ outside the discs, by Lemma~\ref{lem:laurent} the
clockwise integrals equal minus the sums of the residues of $\Delta^{\mathrm{mod}}(s)/(s-z)$ at $s=0$ and $s=4c$, that is,
$\mathfrak A_0/z+\mathfrak A_r/(z-4c)+\mathfrak B_r/(z-4c)^2$. For the coefficients at infinity,
$\mathsf R_k=-\frac1{2\pi i}\int_{\Sigma_{\mathsf R}}\mathsf R_-(J_{\mathsf R}-I)s^{k-1}\diff s$ with the same splitting; on the far arcs
$|s|^{k-1}\le Cn^p$ is absorbed by $e^{-n/C}$, and $-\frac1{2\pi i}\oint_{\mathrm{cw}}\Delta^{\mathrm{mod}}(s)s^{k-1}\diff s$ is the sum of the
residues of $\Delta^{\mathrm{mod}}(s)s^{k-1}$, which is $\mathfrak A_0+\mathfrak A_r$ for $k=1$ and $4c\mathfrak A_r+\mathfrak B_r$ for $k=2$
(the simple pole at $0$ is removed by the factor $s$).
\end{proof}

\subsection{\texorpdfstring{Proof of Theorem~\ref{thm:onen}}{Proof of Theorem 1.4}}\label{sec:second-proof}

By Proposition~\ref{prop:frame}(viii) and Lemma~\ref{lem:invariance}, $n^{2p}b_n^2=\mathsf b(Q)$ and $n^pa_n=\mathsf a(Q)$. For $|z|$
large, $Q=\mathsf RP^\infty_0$, so $Q_1=P_1+\mathsf R_1$ and $Q_2=P_2+\mathsf R_1P_1+\mathsf R_2$ with $P_1,P_2$ of~\eqref{eq:P1P2}. Write
$\mathsf R_1=\mathcal A/n+O(n^{-2})$ and $\mathsf R_2=\mathcal B/n+O(n^{-2})$, with $\mathcal A:=\mathfrak A_0+\mathfrak A_r$ and
$\mathcal B:=4c\mathfrak A_r+\mathfrak B_r$ (Proposition~\ref{prop:second}).

\emph{The coefficient $b_n$.} We have
\[
  n^{2p}b_n^2=\Bigl(ic+\frac{\mathcal A_{12}}n\Bigr)\Bigl(-ic+\frac{\mathcal A_{21}}n\Bigr)+O(n^{-2})
  =c^2\Bigl(1+\frac{i(\mathcal A_{21}-\mathcal A_{12})}{cn}\Bigr)+O(n^{-2}).
\]
The off-diagonal entries of $\mathfrak A_0$ are equal, so by~\eqref{eq:Ar-combinations}
$i(\mathcal A_{21}-\mathcal A_{12})/c=i\cdot ipc(1+\alpha)/c=-p(1+\alpha)=-(s-1)$. Since $b_n>0$,
$n^pb_n=c(1-(s-1)/(2n)+O(n^{-2}))$.

\emph{The coefficient $a_n$.} Since $(Q_1)_{12}=ic+O(1/n)$ and $\mathsf a$ is smooth near $(P_1,P_2)$, expanding
$\mathsf a(Q)=(Q_2)_{12}/(Q_1)_{12}-(Q_1)_{22}$ to first order, with $(\mathsf R_1P_1)_{12}=ic(\mathsf R_1)_{11}$, gives
\[
  n^pa_n=2c+\frac1n\Bigl(\mathcal A_{11}-\mathcal A_{22}+2i\mathcal A_{12}-\frac ic\mathcal B_{12}\Bigr)+O(n^{-2}).
\]
For $\mathfrak A_0=k\bigl(\begin{smallmatrix}1&i\\i&-1\end{smallmatrix}\bigr)$ the bracket vanishes ($k(1+1+2i\cdot i)=0$), and
$\mathfrak A_0$ does not enter $\mathcal B$. The Airy part of the bracket is
$(\mathfrak A_r)_{11}-(\mathfrak A_r)_{22}+2i(\mathfrak A_r)_{12}-4i(\mathfrak A_r)_{12}-\frac ic(\mathfrak B_r)_{12}=-pc(\alpha+2)$
by~\eqref{eq:Ar-combinations}; the terms $12\alpha^2p$ of $\mathfrak A_r$ (Lemma~\ref{lem:residues-full}) cancel between the
contributions of $\mathsf R_1$ and $\mathsf R_2$. Since $p(\alpha+2)=s-1+p$, $n^pa_n=2c(1-(s-1+p)/(2n)+O(n^{-2}))$.

\emph{Consequences.} Dividing, $a_n/(2b_n)=(1-\frac{s-1+p}{2n})(1-\frac{s-1}{2n})^{-1}+O(n^{-2})=1-\frac p{2n}+O(n^{-2})$.
Finally $c(n+\frac{s-1}{2p})^{-p}=cn^{-p}(1-\frac{s-1}{2n}+O(n^{-2}))$, which gives the last form of $b_n$.
\qed

\begin{remark}[The hard edge at order $1/n$; the case $p=2$]\label{rem:second}
The hard edge contributes nothing at order $1/n$: its residue $\mathfrak A_0$ is a multiple of
$\mathsf Te_1e_2^{\mathsf T}\mathsf T^{-1}$, because only the entry $a^2Y_{12}$ in~\eqref{eq:conj-T} has a pole at $0$, and this
direction is annihilated by both functionals. Hence the $1/n$ terms of $b_n$ and $a_n$ are those of the Airy correction
at the soft edge, and of these only the residue combinations~\eqref{eq:Ar-combinations} enter; the coefficient of $b_n$ comes
from the simple pole alone. At $p=2$ the band factor is exactly constant (the remark after Lemma~\ref{lem:bandfactor}), so
$\Delta_0=\Delta_0^{\mathrm{mod}}$ and $\Delta_r=\Delta_r^{\mathrm{mod}}$ exactly, and the closed forms of
Remark~\ref{rem:lommel} for $(p,s)=(2,2)$ give the values of Theorem~\ref{thm:onen} there.

\emph{Derivation of the closed form of Remark~\ref{rem:lommel}.} Put
$\beta_{n+1}:=[2\sqrt{(n+\nu)(n+\nu+1)}]^{-1}$. Van Assche \cite[\S3]{VA} records that the orthonormal Lommel
polynomials satisfy $xp_n=\beta_{n+1}p_{n+1}+\beta_np_{n-1}$. For
$\nu=\frac32$ the orthogonality measure $\mu$ is identified as follows. By the Bessel recurrence
\cite[(10.6.1)]{DLMF}, $\mathsf u_n:=\sqrt{n+\nu}\,J_{\nu+n}(1/x)$ solves $x\mathsf u_n=\beta_{n+1}\mathsf
u_{n+1}+\beta_n\mathsf u_{n-1}$ for $n\ge0$ (with $\beta_0:=[2\sqrt{\nu(\nu-1)}]^{-1}$), and it is square summable by
\cite[(10.2.2)]{DLMF}. For $x\notin\R$ the resolvent vector $(\mathcal J-x)^{-1}e_0$ of the Jacobi matrix $\mathcal J$
solves the same recurrence for $n\ge1$ and is square summable, and such solutions form a line, since the Wronskian
$\beta_n(\mathsf u_n\tilde{\mathsf u}_{n-1}-\mathsf u_{n-1}\tilde{\mathsf u}_n)$ of two of them is constant and tends to $0$. Hence
$\int(t-x)^{-1}\diff\mu(t)$, the first entry of the resolvent, equals $-\mathsf u_0/(\beta_0\mathsf u_{-1})=-2\nu
J_\nu(1/x)/J_{\nu-1}(1/x)$. For $\nu=\frac32$ this is $-3(z^{-1}-\cot z)$ with $z=1/x$ \cite[(10.16.1),
(10.6.1)]{DLMF}, and the partial fractions of $\cot$ \cite[(4.22.3)]{DLMF} show that $\mu$ has mass $3/(\pi k)^2$ at
$\pm1/(\pi k)$, $k\ge1$. Its image under $t\mapsto\pi^2t^2$ is $(6/\pi^2)\sigma_{2,2}$. The quadratic transformation
(the monic orthogonal polynomials of $\mu$ of degree $2n$ are $q_n(t^2)$, where $q_n$ are those of the image of $\mu$
under $t\mapsto t^2$) gives $a_n=\pi^2(\beta_{2n}^2+\beta_{2n+1}^2)$ and $b_n=\pi^2\beta_{2n-1}\beta_{2n}$ for
$\sigma_{2,2}$, which are the formulas of Remark~\ref{rem:lommel}.
\end{remark}
